% Part: incompleteness
% Chapter: theories-computability
% Section: first-incompleteness

\documentclass[../../../include/open-logic-section]{subfiles}

\begin{document}

\olfileid{inc}{tcp}{inc}

\olsection{$\Th{Q}$ चा संपूर्ण, सुसंगत, !!^{axiomatizable}
  विस्तार नाही}

\begin{thm}
\ollabel{thm:first-incompleteness}
$\Th{Q}$ चा संपूर्ण, सुसंगत आणि !!{axiomatizable} विस्तार
अस्तित्वात नाही.
\end{thm}

\begin{proof}
$\Th{Q}$ चा कोणताही सुसंगत, निर्णेय विस्तार नसतो, हे
आपल्याला आधीच माहीत आहे. पण $\Th{T}$ संपूर्ण आणि
!!{axiomatizable} असेल, तर ती निर्णेय असते.
\end{proof}

\begin{explain}
हे प्रमेय गोडेल यांच्या १९३१ मधील पहिल्या अपूर्णता प्रमेयाच्या
मूळ मांडणीपासून फारसे दूर नाही. अधिक आधुनिक पारिभाषिक शब्द
सोडले, तर मुख्य फरक असे आहेत: गोडेल यांनी ``सुसंगत''ऐवजी
``$\omega$-सुसंगत'' ही अट वापरली; तसेच संगणनक्षमतेची आकारिक
संकल्पना तेव्हा प्रस्थापित झाली नव्हती, म्हणून ``!!{axiomatizable}''
हे सर्वसाधारणपणे म्हणता येत नव्हते. (पुढील दशकभरात गोडेल
यांच्यासह अनेकांनी संगणनक्षमतेची आकारिक प्रतिमाने विकसित केली;
गोडेल प्रकारची घटना कोणत्या उपपत्तींना लागू होते हे नेमके
सांगणे, हा त्यामागील एक मोठा उद्देश होता.)

संपूर्णता, सुसंगतता आणि !!{axiomatizability} हे तिन्ही गुणधर्म
एकाच वेळी मिळू शकत नाहीत, असे प्रमेय सांगते. मात्र यांपैकी
कोणताही एक सोडल्यास उरलेले दोन मिळू शकतात: $\Th{Q}$ सुसंगत
आणि संगणनक्षम स्वयंसिद्धकांनी निरूपित आहे, पण संपूर्ण नाही;
विसंगत उपपत्ती संपूर्ण असून संगणनक्षम स्वयंसिद्धकांनी
(उदाहरणार्थ, $\{ \eq/[0][0] \}$ या संचाने) निरूपित होते,
पण सुसंगत नाही; आणि अंकगणितातील सर्व सत्य वाक्यांचा संच
संपूर्ण व सुसंगत आहे, पण संगणनक्षम स्वयंसिद्धकांनी
निरूपित होत नाही.
\end{explain}
\end{document}
